\section{Validation of Feasibility}

Much of the proposed device is already well supported by existing literature. The wristwatch form factor and photoplethysmography (PPG) sensing modality are widely adopted in consumer and research-grade wearables, and heart rate (HR) estimation from wrist PPG is a relatively mature problem, with numerous validated algorithms achieving high accuracy even during motion.

Respiration rate (RR) estimation from wrist PPG, however, remains an active area of research, largely because the respiratory-induced modulations of the PPG waveform are considerably weaker and more susceptible to corruption than the pulsatile component used for HR estimation. There is currently no consensus on the optimal signal processing pipeline for extracting RR from wrist PPG, and reliable RR estimation during motion, such as exercise, is essentially unsolved.

Given that the proposed device relies on RR as a key physiological input for estimating breathlessness, it is necessary to establish whether a reasonable RR estimate can be obtained from wrist PPG under realistic conditions before committing to this sensing approach. To this end, this section presents a validation analysis using raw PPG data from two randomly selected participants performing sitting and walking activities, drawn from the publicly available PPG-DaLiA dataset \cite{reissPPGDaLiA2019}. The two participants, referred to throughout by their dataset identifiers S7 and S8, also offer a useful contrast in self-reported baseline fitness: S7 is a 21-year-old female who reported rarely engaging in sport activity, while S8 is a 43-year-old male who reported exercising several times per week. This analysis does not attempt to solve the broader problem of reliable RR estimation from wrist PPG during exercise, but instead aims to validate that a reasonable estimate can be obtained when compared against a ground truth respiration measurement. If reliable RR data can indeed be obtained in this way, then coupled with HR data, it would be recommended to conduct a further study into whether these two parameters alone form a sufficient basis to classify a patient's breathlessness severity, compared against their own self-assessment. Furthermore, despite this not addressing RR estimation from wrist PPG during motion itself, this is the intended use of the device: gathering this data during rest periods when the wearer is still, such as breaks during sports or rests between weightlifting exercises, to guide the wearer through breathing exercises that aid their recovery.

\subsection{Ground Truth Derivation and the Effect of Motion}

For each participant's sitting and walking sections, a respiration rate ground truth was derived from the chest-mounted RespiBAN unit's respiration (impedance pneumography) signal, rather than being taken from a single breath-detection algorithm. Three independent breath-detection methods -- time-domain peak detection, zero-crossing counting, and spectral dominant-frequency estimation -- were each run over 8~s windows (2~s shift) of the band-pass filtered (0.10--0.50~Hz) chest respiration signal, and a window's rate was only accepted once at least two of the three methods agreed to within 6 breaths/min; disagreeing or low-quality windows were rejected rather than reported. For the sitting sections, shown in \figref{fig:gt-sitting}, this ensemble produced an accepted rate for 97\% (S7) and 89\% (S8) of windows, with a smoothly varying respiration rate (mean $19.7 \pm 1.4$ and $18.9 \pm 1.6$ breaths/min respectively) consistent with expected resting adult respiration rates -- giving confidence that this is a reasonable ground truth to validate the wrist PPG-derived estimate against.

\begin{figure}[t]
    \centering
    \subfloat[S7]{\includegraphics[width=0.95\columnwidth]{S7_resp_rate_sitting.png}}\\
    \subfloat[S8]{\includegraphics[width=0.95\columnwidth]{S8_resp_rate_sitting.png}}
    \caption{Ground truth respiration rate estimation for the sitting sections of (a) S7 and (b) S8: raw chest RespiBAN respiration signal (top) and the resulting fused rate at each accepted window, with a smoothed trend (bottom). The high proportion of accepted windows (97\% S7, 89\% S8) and low estimate variance ($\pm1.4$--$1.6$ breaths/min) support the use of this ensemble as ground truth for validating the wrist-PPG-derived estimate.}
    \label{fig:gt-sitting}
\end{figure}

The walking sections were considerably less reliable: the accepted-window rate dropped to 39\% (S7) and 46\% (S8), and the windows that were still accepted were far noisier (standard deviation rising from $\sim$1.5 to $\sim$5 breaths/min). \figref{fig:ppg-noise} illustrates the source of this degradation at the wrist: the raw wrist PPG signal, which shows a clean, near-periodic pulse train of amplitude $\sim$20--100 (a.u.) during sitting, is dominated by large, irregular motion artefacts (peaks exceeding $\pm 400$) throughout the walking section, with no consistently identifiable pulse structure. Because the chest strap is more rigidly coupled to the torso than a wristband is to the wrist, its signal degrades less severely under motion, but even this reference becomes considerably less trustworthy while walking. Consistent with the literature discussed above, RR estimation from the wrist PPG during walking was therefore not pursued further in this validation; the remainder of this section focuses on the sitting data, where both the ground truth and the wrist-PPG-derived estimate are reasonably trustworthy.

\begin{figure}[t]
    \centering
    \subfloat[Sitting]{\includegraphics[width=0.95\columnwidth]{S7_ppg_signal_sitting.png}}\\
    \subfloat[Walking]{\includegraphics[width=0.95\columnwidth]{S7_ppg_signal_walking.png}}
    \caption{Raw wrist PPG (BVP) signal for participant S7 during (a) sitting and (b) walking, over the same amplitude-comparable time windows. The near-periodic pulse train visible at rest is replaced by large, irregular motion artefacts once walking, with no consistently identifiable pulse structure -- the reason wrist-PPG respiration estimation during motion was not attempted in this validation.}
    \label{fig:ppg-noise}
\end{figure}

\subsection{Respiration Rate Estimation from Wrist PPG (Sitting)}

For the sitting sections, respiration rate was estimated from the wrist PPG signal alone, with no reference to the chest strap, by exploiting three independent ways in which respiration modulates a PPG pulse train: frequency modulation of the pulse-to-pulse interval (respiratory sinus arrhythmia), amplitude modulation of pulse amplitude (respiration-driven changes in stroke volume), and baseline wander of the PPG's DC level (respiration-driven changes in peripheral blood volume). Individual heartbeats were first identified from the wrist PPG using an incremental-merge peak/onset detector, with beat-level rejection of implausible pulse intervals or amplitudes. Each of the three respiration-modulated streams was then reconstructed on a common time grid from these detected beats, band-pass filtered to the breathing band (0.10--0.50~Hz), and converted into a breathing rate every 2~s using peak-interval statistics from a rolling 24~s window; the three streams' rate estimates were fused with a quality-weighted vote, discarding a window's estimate where the streams disagreed excessively.

An initial version of this pipeline reproduced the correct within-window shape of the reference respiration trend but consistently underestimated its level, by around 3 breaths/min for S7 and 4 breaths/min for S8 -- a fairly constant per-participant offset rather than one that varied over the session. This is expected: none of the three respiration-modulated streams carries any absolute information about a person's mean breathing rate, only about its relative fluctuation, so the fused estimate settles at whatever level its own beat-to-beat modulations happen to sit at. The pipeline was therefore revised to include a one-off per-participant DC baseline calibration: a single additive offset ($+2.5$ breaths/min for S7, $+3.6$ breaths/min for S8), fitted once against the reference respiration rate over the session, analogous to a short calibration a real device might perform once against a reference measurement before regular use. This step does not alter the underlying detection method or the breath-to-breath shape it recovers -- only the constant level it is centred on.

\figref{fig:s7-rr} and \figref{fig:s8-rr}, given in the appendix, show the resulting respiration rate for S7 and S8 respectively, compared against the RespiBAN ground truth over the full sitting section. After DC calibration, the mean absolute error against the ground truth fell from 3.1 to 1.8 breaths/min for S7, and from 4.3 to 2.2 breaths/min for S8 (RMSE: 3.8 to 2.3, and 5.0 to 2.8 breaths/min respectively). The estimate does not track every individual fluctuation in the reference trace -- window-to-window correlation remains weak -- reflecting how subtle the respiratory modulation of a wrist PPG pulse is relative to the noise floor discussed in the introduction to this section, but the overall level and slow trend are well matched throughout both six- to ten-minute sessions. As a secondary check, heart rate derived from the same wrist PPG beat detector was compared against the chest-strap reference heart rate over the same windows: for S7 the tracking was close (MAE 2.0~bpm), while for S8 it was periodically corrupted by isolated motion artefacts even during nominal sitting, producing brief overestimation spikes (up to 110~bpm against a true rate near 60~bpm) and a correspondingly higher MAE (9.6~bpm). Heart rate estimation is not the focus of this validation, but this nonetheless supports the premise that pulse-level HR detection from wrist PPG is the more mature, lower-risk half of the sensing pipeline relative to RR.

\subsection{Summary}

This validation demonstrates that a reasonable respiration rate estimate can be derived from wrist PPG data alone, without any additional sensing modality, provided the wearer is reasonably still. It does not establish whether HR and RR alone are sufficient to predict patient-reported dyspnoea with strong accuracy; a further study, using a validated dyspnoea scale as ground truth, would be needed to determine this. It was also found that motion artefacts severely degrade the accuracy of the wrist-PPG RR estimate, a roadblock already well reported in the literature and consistent with the difficulty of RR estimation from wrist PPG during motion discussed at the outset of this section. The methodology used here is furthermore simple and illustrative rather than conclusive: a larger sample, spanning a range of genders, ages and activity levels across a wider range of activities, should be analysed before drawing firm conclusions, but this is outside the scope of the present report. The key takeaway is that signal processing of the raw wrist PPG signal is the most integral and challenging aspect of realising this product. This validation shows that it is achievable when the wearer is reasonably still -- a condition well matched to the product's intended use, in which the recovery phases of exercise (for example, breaks during sports or rests between weightlifting sets) are used to deliver guided breathing exercises tailored to the wearer's own derived heart rate and respiration rate.
